\documentclass[11pt,letterpaper]{article}
\usepackage[margin=1in]{geometry}
\usepackage{booktabs}
\usepackage{hyperref}
\usepackage{enumitem}
\usepackage{titlesec}
\usepackage{amsmath}
\usepackage{xcolor}
\usepackage{microtype}
\usepackage{parskip}

\hypersetup{colorlinks=true, linkcolor=blue, urlcolor=blue, citecolor=blue}

\titlespacing*{\section}{0pt}{10pt}{4pt}
\titlespacing*{\subsection}{0pt}{8pt}{3pt}

\title{\textbf{Milestone Report: DistServe Reimplementation}\\
\large Separating Prefill and Decode for LLM Serving}
\author{Systems for Deep Learning --- Course Project}
\date{March 28, 2026}

\begin{document}
\maketitle
\vspace{-8pt}
\noindent\rule{\linewidth}{0.4pt}

%% ============================================================
\section{Project Goal}
%% ============================================================

The goal of this project is to reimplement the central systems insight of
\textbf{DistServe}~\cite{zhong2024distserve}: separating LLM inference into a
dedicated \emph{prefill pool} and a dedicated \emph{decode pool} eliminates
prefill--decode interference and improves SLO goodput and tail latency under mixed
workloads. We rebuild this in pure Python + PyTorch + HuggingFace (no Ray, no
SwiftTransformer), targeting a 1-node / 2-GPU environment, and evaluate on real GPU
hardware.

%% ============================================================
\section{What Has Been Completed}
%% ============================================================

\subsection{System Architecture}

We implement two serving designs, each runnable in two modes:

\begin{center}
\begin{tabular}{lll}
\toprule
\textbf{Design} & \textbf{Simulator} & \textbf{Real GPU} \\
\midrule
Colocated (baseline)             & \texttt{src/simulator/baseline.py}      & \texttt{src/runtime/baseline\_gpu.py} \\
Disaggregated (DistServe-style)  & \texttt{src/simulator/disaggregated.py} & \texttt{src/runtime/disaggregated\_gpu.py} \\
\bottomrule
\end{tabular}
\end{center}

Supporting modules: \texttt{src/core/\{request,metrics\}.py},
\texttt{src/simulator/\{timing,workload,config\}.py},
\texttt{src/runtime/inference\_core.py},
\texttt{src/stage\_b/\{profile\_sharegpt,fitted\_timing,workload\_sharegpt\}.py}.

\subsection{Real GPU Runtime (Primary Contribution)}

\texttt{src/runtime/inference\_core.py} provides wall-clock timed wrappers around
HuggingFace \texttt{forward} calls with CUDA synchronization:

\begin{itemize}[noitemsep,topsep=2pt]
  \item \textbf{\texttt{timed\_prefill}} --- single-request prefill forward; records
        wall time bracketed by \texttt{torch.cuda.synchronize}.
  \item \textbf{\texttt{timed\_decode\_steps}} --- greedy autoregressive decode loop;
        records per-step wall time.
  \item \textbf{\texttt{time\_transfer}} --- measures the wall time of moving
        \texttt{past\_key\_values} from the prefill GPU to the decode GPU via
        PyTorch \texttt{.to(device)} (PCIe/NVLink device-to-device copy).
  \item \textbf{\texttt{tokenize\_batch}, \texttt{timed\_prefill\_batch},
        \texttt{timed\_decode\_steps\_batch}} --- new batched variants using
        left-padded tokenization; group $B$ requests into a single forward call
        for both prefill and each decode step.
\end{itemize}

\textbf{Colocated GPU runtime} (\texttt{baseline\_gpu.py}): runs prefill then greedy
decode on a single GPU in FIFO order, with configurable \texttt{batch\_size}.

\textbf{Disaggregated GPU runtime} (\texttt{disaggregated\_gpu.py}): prefill on
GPU~A, measured KV transfer via \texttt{time\_transfer}, decode on GPU~B.
Batching moves the full batched KV cache as a single transfer call.
Both GPUs load TinyLlama/TinyLlama-1.1B-Chat-v1.0 (identical weights, separate
device instances), matching the DistServe architecture of role-specific replicas.

\textbf{KV cache transfer implementation:} after prefill we read each transformer
layer's cached keys and values from HuggingFace's \texttt{DynamicCache},
move every tensor with \texttt{.to(decode\_device)}, and rebuild a fresh
\texttt{DynamicCache} on the decode GPU. This is a same-node operation
(both GPUs in the same process). Cross-node handoff is out of scope for this
milestone and will be stated as a limitation in the final report.

\subsection{Stage B: ShareGPT Profiling Pipeline}

\texttt{src/stage\_b/profile\_sharegpt.py} measures actual GPU prefill and per-step
decode latency on ShareGPT prompts, fits a linear prefill model
($a + b \cdot \text{tokens}$) via least squares, and writes
\texttt{results/fitted\_timing.json}. This pipeline is implemented and can be used
to calibrate the simulator with real-hardware timing coefficients.

\subsection{Simulator (Supporting Infrastructure)}

A discrete-event simulator (\texttt{src/simulator/}) models both designs using a
parametric \texttt{TimingModel}. It is used for large-scale sensitivity analysis
(KV transfer cost sweep, capacity split sweep, arrival rate sweep) where real GPU
runs are too slow to sweep exhaustively. Simulator results are not presented as
primary evidence; they serve as a design and exploration tool.

\subsection{Experiment Infrastructure}

\begin{itemize}[noitemsep,topsep=2pt]
  \item \texttt{run\_comparison.py} --- core colocated vs disaggregated simulator run
  \item \texttt{run\_ablations.py} --- KV transfer cost + capacity split sweeps
  \item \texttt{run\_load\_sweep.py} --- simulator goodput vs arrival rate curve
  \item \texttt{run\_workload\_mix.py} --- simulator goodput vs interactive fraction
  \item \texttt{run\_batch\_sweep\_gpu.py} --- \textbf{real GPU} batch-size sweep
        (batch\_size $\in \{1,2,4,8\}$ for both designs)
  \item \texttt{run\_milestone.sh} --- master script; auto-detects GPU count,
        runs all experiments, writes results to \texttt{results/milestone/}
\end{itemize}

All outputs are timestamped CSVs + JSON metadata for reproducibility.

%% ============================================================
\section{Real GPU Evidence}
%% ============================================================

\textbf{Setup:} TinyLlama/TinyLlama-1.1B-Chat-v1.0, 2-GPU node,
30 ShareGPT-backed requests, \texttt{batch\_size}~=~4.
TTFT SLO~=~2.0\,s, TPOT SLO~=~0.05\,s.
Result file: \texttt{results/20260328T000851Z-gpu-comparison-summary.json}.
Per-request data: \texttt{results/20260328T000851Z-gpu-\{colocated,disaggregated\}-requests.csv}
(screenshot attached as evidence).

\begin{center}
\begin{tabular}{lrrr}
\toprule
\textbf{Metric} & \textbf{Colocated} & \textbf{Disaggregated} & \textbf{Change} \\
\midrule
Mean TTFT (s)        & 3.994 & 3.414 & \textbf{--14.5\%} \\
p50 TTFT (s)         & 3.506 & 2.925 & --16.7\% \\
p95 TTFT (s)         & 7.256 & 6.675 & --8.0\% \\
Mean E2E latency (s) & 4.347 & 3.766 & --13.4\% \\
p99 E2E latency (s)  & 8.112 & 7.516 & --7.3\% \\
Throughput (req/s)   & 3.686 & 3.977 & \textbf{+7.9\%} \\
SLO Goodput          & 0.40  & 0.40  & tied \\
\bottomrule
\end{tabular}
\end{center}

\textbf{Interpretation.}
The disaggregated design reduces mean TTFT by 14.5\% and improves throughput by
7.9\% over the colocated baseline on identical requests and hardware. Both designs
achieve the same SLO goodput (0.40) at this batch size and SLO threshold because
the queuing effect at \texttt{batch\_size}~=~4 pushes mean TTFT well above the
2.0\,s threshold for a large fraction of requests in both cases.
The raw latency reduction is real and directionally consistent with the DistServe
claim: separating prefill from decode reduces the time until the first token is
produced. An earlier sequential run (\texttt{batch\_size}~=~1, 15 requests)
showed goodput 0.333~$\to$~0.467 (+40\%) and mean TTFT
2.629~$\to$~2.342\,s (--10.9\%), confirming the direction holds across batch sizes.

\textbf{TPOT.} Mean TPOT is similar across designs (0.0192\,s vs 0.0191\,s),
which is expected: once KV is on the decode GPU, per-step decode time is governed
by the model and memory bandwidth, not by the prefill--decode split.

%% ============================================================
\section{What Remains Before the Final Report}
%% ============================================================

\begin{enumerate}[noitemsep,topsep=3pt]

  \item \textbf{Real GPU batch-size sweep (\texttt{run\_batch\_sweep\_gpu.py}).}
  The script is implemented but not yet run. Will produce measured throughput and
  TTFT for \texttt{batch\_size} $\in \{1,2,4,8\}$ for both designs on real
  hardware, showing how batching interacts with the prefill--decode split.

  \item \textbf{Real GPU arrival rate sweep.}
  The load sweep currently exists only in the simulator. Running it on real
  hardware --- varying arrival rate across both designs and recording goodput and
  tail latency --- is the most important remaining experiment. It would directly
  validate the simulator's load-collapse story with measured numbers.

  \item \textbf{SLO threshold sweep.}
  Fix the workload and vary the TTFT SLO (e.g.\ 1.0\,s, 1.5\,s, 2.0\,s, 3.0\,s)
  to produce a goodput-vs-SLO curve for both designs. This is the canonical
  evaluation plot from the DistServe paper and shows the advantage holds across
  different strictness levels.

  \item \textbf{Direct interference measurement.}
  Run prefill-only vs prefill+decode colocated on the same GPU and directly
  measure the slowdown. This empirically quantifies the interference the
  simulator approximates with a fixed 1.35$\times$ multiplier and removes the
  need for that assumption.

  \item \textbf{Larger GPU runs and longer prompts.}
  Current GPU comparison uses 30 requests truncated to 512 tokens. Scaling to
  100--200 requests with full ShareGPT prompts (1024--2048 tokens) will reduce
  variance, produce stable tail-latency percentiles, and stress the prefill phase
  more — making the KV transfer cost and interference effect more pronounced.

  \item \textbf{Larger model.}
  TinyLlama (1.1B) has short prefill times. Running on a larger model
  (e.g.\ Llama-2-7B if GPU memory permits) would produce longer prefills,
  stronger interference, and a larger disaggregation benefit --- closer to the
  real-world regime the paper targets.

  \item \textbf{KV transfer microbenchmark.}
  Sweep prompt length vs measured \texttt{time\_transfer} output to fit a
  \texttt{transfer\_per\_prompt\_token} coefficient, replacing the current
  synthetic placeholder and closing the loop between the simulator and the GPU
  runtime.

  \item \textbf{Stage B timing calibration.}
  Run \texttt{profile\_sharegpt.py} to produce \texttt{fitted\_timing.json}
  and rerun simulator sweeps with measured timing coefficients, grounding the
  simulator curves in real GPU numbers.

  \item \textbf{Plots and figures.}
  Convert result CSVs into matplotlib figures: goodput vs arrival rate,
  goodput vs SLO threshold, TTFT CDF, batch-size throughput bar chart.
  These will be the primary visual evidence in the final report.

\end{enumerate}

%% ============================================================
\section{On Track / Blockers}
%% ============================================================

\textbf{On track.} The full real-GPU pipeline is working end-to-end: prefill,
KV transfer, and decode are all measured on real hardware. The batching extension
is implemented and tested. The experiment infrastructure is automated and ready
to produce the remaining GPU runs.

\textbf{No hard blockers.} All remaining work is execution and analysis: run more
GPU experiments, calibrate timing, generate figures. No architectural changes are
needed.

\textbf{Scope boundary.} KV transfer is same-node only (device-to-device via
\texttt{.to()}). Multi-node / NCCL-based handoff is out of scope and will be
documented as a limitation.

\textbf{Code structure:}
\begin{itemize}[noitemsep,topsep=2pt]
  \item \texttt{src/runtime/} --- real-GPU runtime with batching support
  \item \texttt{src/simulator/} --- discrete-event simulator (supporting tool)
  \item \texttt{src/stage\_b/} --- ShareGPT profiling and timing fitting
  \item \texttt{src/experiments/} --- five experiment scripts + master runner
  \item \texttt{results/} --- timestamped per-run CSVs and JSON summaries
  \item \texttt{results/milestone/} --- named milestone outputs
  \item \texttt{tests/} --- correctness tests (simulator invariants, timing fit)
\end{itemize}

\vspace{6pt}
\noindent\rule{\linewidth}{0.4pt}
\begin{thebibliography}{1}
\bibitem{zhong2024distserve}
Yinmin Zhong, Shengyu Liu, Junda Chen, Jianbo Hu, Yibo Zhu, Xuanzhe Liu,
Xin Jin, and Hao Zhang.
\newblock DistServe: Disaggregating Prefill and Decoding for Goodput-Optimized
Large Language Model Serving.
\newblock In \textit{OSDI}, 2024.
\end{thebibliography}

\end{document}
